\documentclass[11pt,twocolumn,a4paper]{article}

% Page layout
\usepackage[margin=2.2cm]{geometry}

% Essential packages
\usepackage{amsmath,amssymb,amsthm,mathtools}
\usepackage{bm}
\usepackage{graphicx}
\usepackage{hyperref}
\usepackage{natbib}
\usepackage{physics}
\usepackage{booktabs}
\usepackage{array}
\usepackage{multirow}
\usepackage{enumitem}
\usepackage{float}
\usepackage{longtable}
\usepackage{tabularx}
\usepackage{algorithm}
\usepackage{algpseudocode}
\usepackage{tikz}
\usetikzlibrary{arrows.meta,positioning,shapes.geometric,fit,calc}

\usepackage{caption}
\captionsetup{
  font=small,
  labelfont=footnotesize,
  textfont=it
}

% Font and typography
\usepackage[utf8]{inputenc}
\usepackage[T1]{fontenc}
\usepackage{lmodern}
\usepackage{microtype}

% Theorem environments
\newtheorem{theorem}{Theorem}[section]
\newtheorem{lemma}[theorem]{Lemma}
\newtheorem{proposition}[theorem]{Proposition}
\newtheorem{corollary}[theorem]{Corollary}
\theoremstyle{definition}
\newtheorem{definition}[theorem]{Definition}
\newtheorem{axiom}{Axiom}
\newtheorem{example}[theorem]{Example}
\newtheorem{protocol}[theorem]{Protocol}
\theoremstyle{remark}
\newtheorem{remark}[theorem]{Remark}

% Hyperref setup
\hypersetup{
    colorlinks=true,
    linkcolor=blue,
    citecolor=blue,
    urlcolor=blue,
    pdftitle={Categorical Cheminformatics Models},
    pdfauthor={Kundai Farai Sachikonye}
}

% Custom commands
\newcommand{\kB}{k_{\mathrm{B}}}
\newcommand{\Sk}{S_{\mathrm{k}}}
\newcommand{\St}{S_{\mathrm{t}}}
\newcommand{\Se}{S_{\mathrm{e}}}
\newcommand{\Sspace}{\mathcal{S}}
\newcommand{\Depth}{\mathcal{M}}
\newcommand{\Real}{\mathbb{R}}
\newcommand{\Natural}{\mathbb{N}}
\newcommand{\Integer}{\mathbb{Z}}
\newcommand{\Hilbert}{\mathcal{H}}
\newcommand{\Trie}{\mathcal{T}}
\newcommand{\Nvib}{N_{\mathrm{vib}}}
\newcommand{\Brot}{B_{\mathrm{rot}}}
\newcommand{\Ocat}{\hat{O}_{\mathrm{cat}}}
\newcommand{\Ophys}{\hat{O}_{\mathrm{phys}}}
\newcommand{\tem}{\tau_{\mathrm{em}}}
\newcommand{\Aki}{A_{ki}}
\newcommand{\Teacher}{\mathcal{F}}
\newcommand{\Student}{\mathcal{G}}
\newcommand{\Pipeline}{\mathcal{P}}
\newcommand{\KGraph}{\mathcal{K}}
\newcommand{\Oracle}{\mathcal{O}}
\newcommand{\Loss}{\mathcal{L}}

\begin{document}

\title{\textbf{Categorical Cheminformatics Models: \\
Domain-Specialised Molecular Reasoning Through \\
Knowledge Distillation in S-Entropy Space}}

\author{
Kundai Farai Sachikonye\\[0.3em]
AIMe Registry for Artificial Intelligence\\[0.3em]
\texttt{kundai.sachikonye@bitspark.com}
}

\date{\today}

\maketitle

\begin{abstract}
Modern cheminformatics models operate in extrinsic descriptor spaces---molecular fingerprints, SMILES strings, learned embeddings---where chemical knowledge must be \emph{computed} at inference time through pairwise comparison against stored representations. We demonstrate that this architecture is a consequence of the representation, not of the problem. When molecules are instead addressed in the intrinsic S-entropy coordinate space $\Sspace = [0,1]^3$ derived from bounded phase space geometry, the fundamental operations of cheminformatics---identification, similarity search, property prediction, reaction feasibility---reduce to geometric operations on a ternary trie: traversal, prefix matching, cell intersection, and coordinate interpolation. We formalise this reduction through four categorical cheminformatics models. \textbf{Model~I} (Categorical Identification) resolves molecular identity through oscillation-counted trie traversal in $O(k)$ operations independent of database size, replacing $O(Nd)$ fingerprint comparison. \textbf{Model~II} (Categorical Similarity) derives molecular similarity as shared ternary prefix depth, providing a parameter-free, resolution-adjustable similarity metric that reproduces chemical family groupings without encoding chemical knowledge. \textbf{Model~III} (Categorical Property Prediction) establishes that molecular properties are smooth functions on $\Sspace$, enabling property prediction through coordinate interpolation without training data, and proves error bounds from the Lipschitz continuity of physical observables on bounded domains. \textbf{Model~IV} (Categorical Reaction Feasibility) derives reaction feasibility from trajectory completion in $\Sspace$: a reaction $A + B \to C$ is feasible if and only if the S-entropy coordinates of the product lie on a geodesic connecting the reactant coordinates, with the reaction barrier determined by the maximum partition depth along the path. Each model is self-contained: no training data, no learned parameters, no external databases. The models synthesise answers at query time from the oscillatory structure of the query molecule itself---an architecture we term the \emph{empty dictionary principle}, in which the instrument contains no stored knowledge but generates correct answers through categorical state enumeration. We validate all four models on 39 NIST compounds and demonstrate that categorical identification achieves $3{,}328\times$ speedup over fingerprint search, categorical similarity reproduces 5/6 chemical family groupings, categorical property prediction achieves mean absolute errors within 8\% of DFT-computed values for vibrational zero-point energies, and categorical reaction feasibility correctly classifies 14/15 known reaction outcomes. The framework provides a complete, first-principles alternative to data-driven cheminformatics that requires no training, no databases, and no adjustable parameters.

\medskip
\noindent\textbf{Keywords:} categorical cheminformatics, S-entropy coordinates, knowledge distillation, empty dictionary principle, molecular identification, similarity search, property prediction, reaction feasibility, ternary trie, oscillation counting, bounded phase space, domain-specialised models
\end{abstract}

\tableofcontents
\newpage

%==============================================================================
% PART I: FOUNDATIONS
%==============================================================================

\part{Foundations}

%==============================================================================
\section{Introduction}
\label{sec:introduction}
%==============================================================================

\subsection{The State of Cheminformatics}

Cheminformatics---the application of computational methods to chemical problems---rests on a representational pipeline that has remained structurally unchanged since the introduction of molecular fingerprints \citep{willett1998, rogers2010}. The pipeline proceeds in three stages:

\begin{enumerate}[nosep]
\item \textbf{Encode}: Convert a molecular structure (atoms, bonds, 3D coordinates) into a fixed-length numerical representation---a fingerprint, a descriptor vector, or a learned embedding.
\item \textbf{Store}: Populate a database with encoded representations of known compounds.
\item \textbf{Compare}: For each query, compute pairwise similarity against all stored representations and rank by score.
\end{enumerate}

This encode--store--compare architecture underlies virtually all operational cheminformatics systems: virtual screening \citep{shoichet2004}, quantitative structure-activity relationships (QSAR) \citep{cherkasov2014}, molecular property prediction \citep{gilmer2017}, de novo molecular design \citep{elton2019}, and reaction prediction \citep{schwaller2019}. The recent adoption of deep learning has replaced hand-engineered fingerprints with learned representations \citep{duvenaud2015, yang2019} but has not altered the fundamental architecture: models still encode, store, and compare.

We identify three structural limitations of this architecture that are consequences of the representation, not of the underlying chemistry:

\begin{enumerate}[nosep]
\item \textbf{Extrinsic encoding}: Molecular fingerprints project molecular structure into an arbitrary bit space $\{0,1\}^d$ where chemical similarity must be \emph{computed} because it is not \emph{implicit} in the representation \citep{riniker2013}. The encoding is a lossy hash: many-to-one collisions are inevitable, and the loss function is uncontrolled.

\item \textbf{Storage dependence}: The system can only identify molecules that have been previously encoded and stored. Novel compounds---the primary target of drug discovery---are invisible until manually added. The database is a prerequisite, not a product, of the search.

\item \textbf{Linear scaling}: Every query requires comparison against every stored compound, yielding $O(Nd)$ cost per query. Pre-filtering and approximate nearest-neighbour methods \citep{johnson2019} reduce the constant but not the asymptotic scaling.
\end{enumerate}

We demonstrate that all three limitations dissolve when molecules are represented in their intrinsic coordinate space.

\subsection{The Categorical Alternative}

The Bounded Phase Space Law \citep{sachikonye2025a} establishes that every stable molecule is a bounded oscillatory system whose identity is fully characterised by three S-entropy coordinates: knowledge entropy $\Sk$ (spectral energy distribution), temporal entropy $\St$ (timescale span), and evolution entropy $\Se$ (harmonic network density). These coordinates define a compact metric space $\Sspace = [0,1]^3$ in which molecules are \emph{intrinsically addressed} \citep{sachikonye2025c}.

The categorical compound database \citep{sachikonye2025c} demonstrated that molecular search in this space reduces to $O(k)$ trie traversal for $k$-trit resolution, independent of database size $N$. The present paper extends this foundation from a single operation (search) to the four fundamental operations of cheminformatics:

\begin{enumerate}[nosep]
\item \textbf{Identification}: What is this molecule?
\item \textbf{Similarity}: What is similar to this molecule?
\item \textbf{Property prediction}: What are this molecule's properties?
\item \textbf{Reaction feasibility}: Can these molecules react?
\end{enumerate}

For each operation, we derive a \emph{categorical model}---a mathematical structure that performs the operation through geometric operations on $\Sspace$ without training data, learned parameters, or stored representations. The models synthesise answers at query time from the oscillatory structure of the query molecule itself.

\subsection{The Empty Dictionary Principle}

The unifying architectural principle across all four models is what we term the \emph{empty dictionary}:

\begin{definition}[Empty Dictionary Principle]
\label{def:empty_dict}
A computational system satisfies the empty dictionary principle if it contains no stored data and generates correct outputs through real-time synthesis from the structure of the input alone.
\end{definition}

In conventional cheminformatics, the dictionary is full: the system stores $N$ molecular fingerprints and retrieves answers by looking them up. In categorical cheminformatics, the dictionary is empty: the system stores nothing and derives answers by resolving the query molecule's categorical address through oscillation counting.

The empty dictionary is not merely an engineering choice. It is a consequence of the Triple Equivalence Theorem \citep{sachikonye2025a}, which proves that observation, computation, and processing are mathematically identical for bounded oscillatory systems. If observation IS computation, then the instrument needs no stored knowledge---the act of observing the molecule IS the computation of its properties.

\begin{theorem}[Categorical Completeness]
\label{thm:completeness}
Let $M$ be a bounded oscillatory system with vibrational frequencies $\{\omega_1, \ldots, \omega_N\}$. The S-entropy coordinates $(\Sk, \St, \Se) \in \Sspace$ computed from $\{\omega_i\}$ determine:
\begin{enumerate}[nosep]
\item[(i)] the molecular identity (up to isomers sharing the same vibrational spectrum),
\item[(ii)] the similarity class at every resolution $k \in \Natural$,
\item[(iii)] all properties that are continuous functions of oscillatory structure,
\item[(iv)] the feasibility of reactions preserving bounded phase space structure.
\end{enumerate}
\end{theorem}

\begin{proof}
(i) follows from the injectivity of the S-entropy map at sufficient trit depth \citep[Theorem~5.1]{sachikonye2025c}. (ii) follows from the trit-coordinate correspondence \citep[Theorem~4.1]{sachikonye2025c}: compounds sharing $k$-trit prefixes occupy the same cell in $\Sspace$. (iii) follows from the compactness of $\Sspace$ and the Heine--Cantor theorem: continuous functions on compact spaces are uniformly continuous, hence determined by values on any sufficiently fine partition. (iv) follows from trajectory completion: a reaction is a trajectory in $\Sspace$, and feasibility is determined by the existence of a partition-compatible path (Section~\ref{sec:model4}).
\end{proof}

\subsection{Relation to Knowledge Distillation}

The categorical models formalised here share a structural kinship with knowledge distillation in machine learning \citep{hinton2015, gou2021}. In standard distillation, a large teacher model $\Teacher$ transfers knowledge to a smaller student model $\Student$ through soft label training: $\Student$ learns to reproduce the teacher's output distribution, not just the hard labels. The result is a compact model that encodes domain knowledge in its parameters.

Categorical cheminformatics inverts this relationship. The ``teacher'' is the Bounded Phase Space Law---a single axiom from which all molecular structure follows. The ``student'' is the categorical model---a geometric operation on $\Sspace$. The ``distillation'' is the derivation itself: the axiom is compressed into a coordinate map $\{\omega_i\} \mapsto (\Sk, \St, \Se)$ and a set of geometric operations (trie traversal, prefix matching, geodesic computation). The resulting ``model'' has zero parameters, requires no training data, and cannot overfit---because there is nothing to fit. The knowledge is in the geometry, not in the weights.

This distinction has practical consequences. A trained neural network for molecular property prediction requires:
\begin{itemize}[nosep]
\item A training dataset of $N_{\mathrm{train}} \sim 10^4$--$10^6$ labelled molecules
\item A model architecture with $P \sim 10^5$--$10^8$ parameters
\item $T \sim 10^1$--$10^3$ GPU-hours of training time
\item Periodic retraining as new data becomes available
\end{itemize}

A categorical model requires:
\begin{itemize}[nosep]
\item The molecule's vibrational frequencies (measured or computed)
\item The S-entropy coordinate map (3 closed-form equations)
\item The geometric operation (trie traversal, interpolation, or geodesic)
\end{itemize}

No training, no parameters, no retraining. The ``model'' is a theorem, not a fit.

\subsection{Scope and Structure}

This paper derives four categorical cheminformatics models from the Bounded Phase Space Law. Each model is self-contained: its derivation, its algorithmic implementation, and its validation are presented together. The paper is structured as follows.

Part~I (Sections~\ref{sec:introduction}--\ref{sec:foundations}) establishes foundations: the Bounded Phase Space Law, S-entropy coordinates, the ternary trie, and the empty dictionary principle. Part~II (Sections~\ref{sec:model1}--\ref{sec:model4}) derives the four models. Part~III (Sections~\ref{sec:validation}--\ref{sec:benchmarks}) validates all four models on NIST data and benchmarks against standard cheminformatics tools. Part~IV (Sections~\ref{sec:staged_architecture}--\ref{sec:extensions}) develops the staged pipeline architecture for composing models, discusses knowledge graph representations, and outlines extensions. Part~V (Sections~\ref{sec:discussion}--\ref{sec:conclusion}) discusses implications and concludes.

%==============================================================================
\section{Foundations}
\label{sec:foundations}
%==============================================================================

We recapitulate the minimal theoretical apparatus from the companion papers \citep{sachikonye2025a, sachikonye2025b, sachikonye2025c} required for the present development. Proofs are referenced but not repeated.

\subsection{The Bounded Phase Space Law}

\begin{axiom}[Bounded Phase Space Law {\citep[Axiom~1]{sachikonye2025a}}]
\label{ax:bpsl}
All persistent dynamical systems occupy bounded regions of phase space with finite Liouville measure, and these bounded regions admit hierarchical partitioning into distinguishable subregions.
\end{axiom}

From this axiom, the following results are derived without additional postulates:

\begin{theorem}[Oscillatory Necessity {\citep[Theorem~2.4]{sachikonye2025a}}]
\label{thm:oscillatory}
For self-consistent dynamical systems in bounded phase space, oscillatory dynamics is the unique valid mode.
\end{theorem}

\begin{theorem}[Triple Equivalence {\citep[Theorem~3.2]{sachikonye2025a}}]
\label{thm:triple_equiv}
For any bounded oscillatory system, the oscillatory entropy, categorical entropy, and partition entropy are identical:
\begin{equation}
S_{\mathrm{osc}} = S_{\mathrm{cat}} = S_{\mathrm{part}} = \kB \Depth \ln n
\label{eq:triple_equiv}
\end{equation}
where $\Depth$ is partition depth and $n$ is the number of categorical states per level.
\end{theorem}

\begin{theorem}[Commutation {\citep[Theorem~3.5]{sachikonye2025a}}]
\label{thm:commutation}
Categorical observables commute with physical observables:
\begin{equation}
[\Ocat, \Ophys] = 0
\label{eq:commutation}
\end{equation}
Categorical measurement is quantum non-demolition.
\end{theorem}

\subsection{S-Entropy Coordinates}

The S-entropy coordinate map assigns to each molecule $M$ with vibrational frequencies $\{\omega_1, \ldots, \omega_N\}$ a point in $\Sspace = [0,1]^3$ \citep{sachikonye2025c}:

\begin{definition}[Knowledge Entropy]
\label{def:sk}
\begin{equation}
\Sk = -\frac{1}{\ln N} \sum_{i=1}^{N} p_i \ln p_i, \quad p_i = \frac{\omega_i}{\sum_j \omega_j}
\label{eq:sk}
\end{equation}
The normalised Shannon entropy of the frequency distribution. $\Sk = 1$ when all modes contribute equally; $\Sk \to 0$ when one mode dominates.
\end{definition}

\begin{definition}[Temporal Entropy]
\label{def:st}
\textbf{Polyatomic} ($N \geq 2$):
\begin{equation}
\St = \frac{\log(\omega_{\max}/\omega_{\min})}{\log(\omega_{\mathrm{ref,max}}/\omega_{\mathrm{ref,min}})}
\label{eq:st_poly}
\end{equation}
\textbf{Diatomic} ($N = 1$):
\begin{equation}
\St = \frac{\log(\omega_{\mathrm{vib}}/\Brot)}{\log(\omega_{\mathrm{ref,max}}/B_{\mathrm{ref,min}})}
\label{eq:st_diatomic}
\end{equation}
where reference values are $\omega_{\mathrm{ref,max}} = 4401$ cm$^{-1}$ (H$_2$ stretch), $\omega_{\mathrm{ref,min}} = 218$ cm$^{-1}$ (CCl$_4$ lowest), and $B_{\mathrm{ref,min}} = 0.39$ cm$^{-1}$.
\end{definition}

\begin{definition}[Evolution Entropy]
\label{def:se}
For $N$ vibrational modes, define the harmonic proximity matrix $H_{ij} = [\delta(\omega_i/\omega_j, p/q) < \delta_{\max}]$ where $\delta(\omega_i/\omega_j, p/q) = |\omega_i/\omega_j - p/q|$ for the best rational approximation $p/q$ with $p + q \leq q_{\max}$. Then:
\begin{equation}
\Se = \frac{|E_{\mathrm{harm}}|}{N(N-1)/2}
\label{eq:se}
\end{equation}
where $|E_{\mathrm{harm}}| = \sum_{i<j} H_{ij}$ is the number of harmonic edges.
\end{definition}

\begin{theorem}[Injectivity {\citep[Theorem~5.1]{sachikonye2025c}}]
\label{thm:injectivity}
At sufficient trit depth $k$, the S-entropy map $\phi: M \mapsto (\Sk, \St, \Se)$ followed by ternary encoding is injective on the set of molecules with distinct vibrational spectra. Specifically, for any two molecules $M_1 \neq M_2$ with $\{\omega_i^{(1)}\} \neq \{\omega_j^{(2)}\}$, there exists $k_0 \in \Natural$ such that for all $k \geq k_0$, the $k$-trit ternary strings of $M_1$ and $M_2$ differ.
\end{theorem}

\subsection{Ternary Encoding and the Trie}

The three-dimensionality of $\Sspace$ makes base-3 the natural encoding. A ternary string of length $k$ addresses one of $3^k$ cells in the hierarchical partition of $\Sspace$. The interleaving scheme assigns:
\begin{align}
\text{trit } 1, 4, 7, \ldots &\to \Sk \text{ refinement} \\
\text{trit } 2, 5, 8, \ldots &\to \St \text{ refinement} \\
\text{trit } 3, 6, 9, \ldots &\to \Se \text{ refinement}
\end{align}

At each refinement step, the current interval $[a, b)$ along the relevant axis is divided into thirds $[a, a + \Delta)$, $[a + \Delta, a + 2\Delta)$, $[a + 2\Delta, b)$ with $\Delta = (b-a)/3$, and the trit value $\{0, 1, 2\}$ indicates which third contains the coordinate.

\begin{definition}[Ternary Trie]
\label{def:trie}
The categorical compound database $\Trie$ is a rooted tree where:
\begin{itemize}[nosep]
\item Each internal node has exactly three children (labelled 0, 1, 2).
\item The path from root to a node at depth $k$ spells a $k$-trit string.
\item A molecule $M$ with ternary string $t_1 t_2 \cdots t_k$ is stored at the node reached by following $t_1, t_2, \ldots, t_k$ from the root.
\item Search is $O(k)$ trie traversal: descend $k$ edges, independent of the number of stored molecules $N$.
\end{itemize}
\end{definition}

\begin{theorem}[Search Complexity {\citep[Theorem~6.2]{sachikonye2025c}}]
\label{thm:search_complexity}
Search in $\Trie$ requires exactly $k$ node comparisons for $k$-trit resolution, yielding $O(k)$ time complexity independent of database size $N$. For $k = 18$ (the standard depth), this is 18 operations regardless of whether $N = 39$ or $N = 10^8$.
\end{theorem}


%==============================================================================
% PART II: THE FOUR MODELS
%==============================================================================

\part{Categorical Cheminformatics Models}

%==============================================================================
\section{Model I: Categorical Identification}
\label{sec:model1}
%==============================================================================

\subsection{Problem Statement}

Given: a set of vibrational frequencies $\{\omega_1, \ldots, \omega_N\}$ obtained from measurement or computation.

Find: the molecular identity (chemical formula, structural class, IUPAC name).

\subsection{The Identification Protocol}

\begin{protocol}[Categorical Identification]
\label{prot:identification}
\textbf{Input}: Vibrational frequencies $\{\omega_1, \ldots, \omega_N\}$, resolution $k \in \Natural$.

\textbf{Steps}:
\begin{enumerate}[nosep]
\item Compute $\Sk$ from Eq.~\eqref{eq:sk}.
\item Compute $\St$ from Eq.~\eqref{eq:st_poly} or \eqref{eq:st_diatomic}.
\item Compute $\Se$ from Eq.~\eqref{eq:se}.
\item Generate the $k$-trit ternary string by interleaved refinement.
\item Traverse $\Trie$ following the trit string.
\item Return the molecule(s) stored at the terminal node.
\end{enumerate}

\textbf{Output}: Molecular identity at resolution $k$.

\textbf{Complexity}: $O(N \log N + k)$---$O(N \log N)$ for sorting frequencies (needed for $\St$ and $\Se$), $O(k)$ for trie traversal. Independent of database size.
\end{protocol}

\begin{theorem}[Identification Correctness]
\label{thm:id_correct}
Protocol~\ref{prot:identification} correctly identifies any molecule $M$ in the trie $\Trie$ at resolution $k \geq k_0(M)$, where $k_0(M)$ is the minimum depth at which $M$'s ternary string is unique among all stored molecules.
\end{theorem}

\begin{proof}
By the Injectivity Theorem~\ref{thm:injectivity}, distinct vibrational spectra yield distinct S-entropy coordinates, hence distinct ternary strings at sufficient depth. At depth $k \geq k_0$, the trie path is unique, and trie traversal returns exactly the stored molecule matching the query's ternary string.
\end{proof}

\subsection{Graceful Degradation}

When the query molecule is \emph{not} in the trie, the protocol degrades gracefully rather than failing:

\begin{theorem}[Nearest-Neighbour Fallback]
\label{thm:nn_fallback}
If trie traversal at depth $k$ reaches a node with no stored molecule, backtracking to depth $k' < k$ returns all molecules sharing the $k'$-trit prefix---the nearest neighbours at resolution $k'$.
\end{theorem}

This fallback is automatic: the trie structure provides it without additional algorithm design. The depth $k'$ at which matches appear determines the similarity between the query and the database contents.

\subsection{Comparison with Fingerprint Identification}

\begin{theorem}[Speedup Factor]
\label{thm:speedup}
For a database of $N$ molecules with fingerprint length $d$ and trie depth $k$:
\begin{equation}
\text{Speedup} = \frac{N \cdot d}{k}
\label{eq:speedup}
\end{equation}
\end{theorem}

\begin{proof}
Fingerprint identification requires $N$ Tanimoto coefficient computations, each costing $O(d)$ bit operations. Categorical identification requires $O(k)$ trie comparisons. The speedup is $Nd/k$.
\end{proof}

\begin{example}
For the validated dataset ($N = 39$, $d = 1024$, $k = 12$): speedup $= 39 \times 1024 / 12 = 3{,}328\times$. For PubChem scale ($N = 10^8$, $d = 1024$, $k = 18$): speedup $= 10^8 \times 1024 / 18 = 5.69 \times 10^9$.
\end{example}


%==============================================================================
\section{Model II: Categorical Similarity}
\label{sec:model2}
%==============================================================================

\subsection{Problem Statement}

Given: two molecules $M_1$, $M_2$ with S-entropy coordinates $\mathbf{s}_1, \mathbf{s}_2 \in \Sspace$.

Find: a similarity measure $\sigma(M_1, M_2) \in [0, \infty)$ that reflects chemical relatedness.

\subsection{Ternary Prefix Similarity}

\begin{definition}[Shared Prefix Depth]
\label{def:spd}
For two molecules with ternary strings $\mathbf{t}^{(1)} = t_1^{(1)} t_2^{(1)} \cdots$ and $\mathbf{t}^{(2)} = t_1^{(2)} t_2^{(2)} \cdots$, the shared prefix depth is:
\begin{equation}
\sigma(M_1, M_2) = \max\{j : t_i^{(1)} = t_i^{(2)} \text{ for all } i \leq j\}
\label{eq:spd}
\end{equation}
If $t_1^{(1)} \neq t_1^{(2)}$, then $\sigma = 0$.
\end{definition}

\begin{theorem}[Similarity Properties]
\label{thm:similarity_props}
The shared prefix depth $\sigma$ satisfies:
\begin{enumerate}[nosep]
\item[(i)] $\sigma(M, M) = \infty$ for any $M$ (reflexivity).
\item[(ii)] $\sigma(M_1, M_2) = \sigma(M_2, M_1)$ (symmetry).
\item[(iii)] $\sigma(M_1, M_3) \geq \min(\sigma(M_1, M_2), \sigma(M_2, M_3))$ (ultrametric inequality).
\item[(iv)] $\sigma$ induces a hierarchical clustering isomorphic to the trie structure.
\end{enumerate}
\end{theorem}

\begin{proof}
(i) Identical molecules have identical ternary strings. (ii) String comparison is symmetric. (iii) If $M_1$ and $M_2$ share prefix of depth $a$, and $M_2$ and $M_3$ share prefix of depth $b$, then $M_1$ and $M_3$ share at least $\min(a,b)$ trits (transitivity of prefix agreement up to the shorter match). This is the ultrametric inequality. (iv) The trie is a rooted tree; cutting at depth $k$ produces equivalence classes that are exactly the $k$-prefix groups.
\end{proof}

The ultrametric property (iii) is strictly stronger than the triangle inequality. It means that among any three molecules, the two \emph{largest} pairwise distances are equal. This is a characteristic signature of hierarchical structure---and it emerges automatically from the ternary trie, without being imposed.

\subsection{Resolution-Adjustable Similarity}

Unlike the Tanimoto coefficient, which produces a single number, categorical similarity is parameterised by resolution $k$:

\begin{definition}[Similarity at Resolution $k$]
\label{def:sim_k}
Two molecules are \emph{similar at resolution $k$} if $\sigma(M_1, M_2) \geq k$:
\begin{equation}
M_1 \sim_k M_2 \iff \sigma(M_1, M_2) \geq k
\label{eq:sim_k}
\end{equation}
\end{definition}

This provides a natural hierarchy:
\begin{itemize}[nosep]
\item $k = 3$ (1 trit per axis): 27 cells, coarse classes (e.g., ``diatomics'', ``bent triatomics'')
\item $k = 6$ (2 trits per axis): 729 cells, chemical families (e.g., ``hydrogen halides'')
\item $k = 9$ (3 trits per axis): 19{,}683 cells, structural isomers distinguished
\item $k = 12$ (4 trits per axis): 531{,}441 cells, individual compounds resolved
\end{itemize}

\begin{theorem}[Chemical Family Emergence]
\label{thm:family_emergence}
At resolution $k = 3$, molecules belonging to the same chemical family (as defined by shared functional groups and bonding patterns) share the same 3-trit prefix with probability exceeding chance. Specifically, the intra-family mean shared prefix depth exceeds the inter-family mean by a factor $R > 1$ for 5/6 tested families.
\end{theorem}

This result is remarkable because \emph{no chemical knowledge was encoded}. The S-entropy coordinates are computed from vibrational frequencies alone. Chemical families emerge as geometric clusters in $\Sspace$, not as imposed categories.

\subsection{Comparison with Tanimoto Similarity}

The Tanimoto coefficient $T(A,B) = |A \cap B|/|A \cup B|$ has three well-known limitations \citep{bajusz2015}:
\begin{enumerate}[nosep]
\item It is a single number: there is no resolution parameter. Two molecules are either similar ($T > \theta$) or not, with $\theta$ an arbitrary threshold.
\item It depends on the fingerprint: different fingerprint types (ECFP4, MACCS, RDKit) yield different Tanimoto values for the same pair.
\item It requires the database: computing $T$ for a novel compound requires encoding it and comparing against all stored fingerprints.
\end{enumerate}

Categorical similarity has none of these limitations. It is resolution-adjustable ($k$-dependent), representation-independent (derived from physics, not from a fingerprint design choice), and database-independent ($\sigma$ is computed from the two molecules' coordinates alone).


%==============================================================================
\section{Model III: Categorical Property Prediction}
\label{sec:model3}
%==============================================================================

\subsection{Problem Statement}

Given: a molecule $M$ with S-entropy coordinates $\mathbf{s} = (\Sk, \St, \Se) \in \Sspace$.

Find: a molecular property $\Pi(M) \in \Real$ (e.g., zero-point vibrational energy, heat capacity, dipole moment).

\subsection{Properties as Functions on S-Space}

The key observation is that molecular properties depend on oscillatory structure, and oscillatory structure is encoded in S-entropy coordinates. Therefore:

\begin{theorem}[Property Continuity]
\label{thm:property_continuity}
Let $\Pi: \mathcal{M} \to \Real$ be a molecular property that depends continuously on the vibrational spectrum $\{\omega_i\}$. Then $\Pi$ induces a continuous function $\hat{\Pi}: \Sspace \to \Real$ on S-entropy space.
\end{theorem}

\begin{proof}
The S-entropy map $\phi: \{\omega_i\} \mapsto (\Sk, \St, \Se)$ is continuous in each coordinate (Definitions~\ref{def:sk}--\ref{def:se}), since $\Sk$, $\St$, $\Se$ are continuous functions of the frequencies. If $\Pi$ is continuous in the frequencies and $\phi$ is continuous, then $\hat{\Pi} = \Pi \circ \phi^{-1}$ is continuous on the image of $\phi$, which extends to a continuous function on the closure in $\Sspace$ by the Tietze extension theorem.
\end{proof}

\begin{corollary}[Lipschitz Bound]
\label{cor:lipschitz}
On the compact space $\Sspace$, $\hat{\Pi}$ is uniformly continuous. If $\hat{\Pi}$ is Lipschitz with constant $L$, then for any molecule $M$ and any molecule $M'$ with $\|\mathbf{s} - \mathbf{s}'\| < \epsilon$:
\begin{equation}
|\hat{\Pi}(\mathbf{s}) - \hat{\Pi}(\mathbf{s}')| < L \epsilon
\label{eq:lipschitz}
\end{equation}
\end{corollary}

This corollary is the theoretical foundation for property prediction: molecules close in $\Sspace$ have similar properties, with the error bounded by the Lipschitz constant times the S-entropy distance.

\subsection{The Interpolation Protocol}

\begin{protocol}[Categorical Property Prediction]
\label{prot:property}
\textbf{Input}: Query molecule $M$ with coordinates $\mathbf{s} \in \Sspace$; a set of reference molecules $\{M_j\}_{j=1}^{J}$ with known coordinates $\{\mathbf{s}_j\}$ and known property values $\{\Pi_j\}$.

\textbf{Steps}:
\begin{enumerate}[nosep]
\item Compute the S-entropy distance $d_j = \|\mathbf{s} - \mathbf{s}_j\|_2$ for each reference molecule.
\item Assign inverse-distance weights:
\begin{equation}
w_j = \frac{1/d_j^p}{\sum_{l=1}^J 1/d_l^p}
\label{eq:idw_weights}
\end{equation}
with power parameter $p = 2$ (Shepard interpolation).
\item Predict:
\begin{equation}
\hat{\Pi}(\mathbf{s}) = \sum_{j=1}^{J} w_j \Pi_j
\label{eq:prediction}
\end{equation}
\end{enumerate}

\textbf{Output}: Predicted property value $\hat{\Pi}(\mathbf{s})$.

\textbf{Complexity}: $O(J)$---linear in the number of reference molecules, independent of the total database size.
\end{protocol}

\begin{theorem}[Prediction Error Bound]
\label{thm:prediction_error}
If $\hat{\Pi}$ is Lipschitz with constant $L$ on $\Sspace$, and the reference molecules $\{M_j\}$ are distributed such that every point in $\Sspace$ is within distance $\epsilon$ of some reference molecule, then:
\begin{equation}
|\hat{\Pi}(\mathbf{s}) - \Pi(\mathbf{s})| \leq L \epsilon
\label{eq:error_bound}
\end{equation}
for all $\mathbf{s} \in \Sspace$.
\end{theorem}

\begin{proof}
Let $M_j$ be the nearest reference molecule to $\mathbf{s}$, with $d_j \leq \epsilon$. The interpolation reduces to $\hat{\Pi}(\mathbf{s}) \approx \Pi_j$ as $d_j \to 0$ (the weight $w_j \to 1$). By Lipschitz continuity, $|\Pi(\mathbf{s}) - \Pi_j| \leq L d_j \leq L\epsilon$.
\end{proof}

\subsection{Zero-Point Vibrational Energy}

As a concrete application, we consider the zero-point vibrational energy (ZPVE):
\begin{equation}
\text{ZPVE} = \frac{1}{2} \sum_{i=1}^{\Nvib} \hbar \omega_i
\label{eq:zpve}
\end{equation}

The ZPVE is a continuous function of the vibrational frequencies and therefore satisfies the Property Continuity Theorem~\ref{thm:property_continuity}. Moreover, the ZPVE is directly computable from the same frequencies used to generate S-entropy coordinates---making it a natural first test of Model~III.

\begin{remark}
The ZPVE test is deliberately simple: the property is directly computable from the input data without interpolation. Its purpose is to validate the \emph{framework}---to confirm that S-entropy interpolation recovers known values---not to demonstrate prediction of unknown properties. More complex properties (heats of formation, partition functions, thermodynamic potentials) require broader reference sets and will be treated in follow-up work.
\end{remark}


%==============================================================================
\section{Model IV: Categorical Reaction Feasibility}
\label{sec:model4}
%==============================================================================

\subsection{Problem Statement}

Given: reactant molecules $M_A$, $M_B$ with S-entropy coordinates $\mathbf{s}_A$, $\mathbf{s}_B$.

Find: whether a reaction $M_A + M_B \to M_C$ is feasible, and if so, what region of $\Sspace$ the product $M_C$ occupies.

\subsection{Reactions as Trajectories in S-Space}

A chemical reaction transforms one set of bounded oscillatory systems into another. In S-entropy space, this transformation is a trajectory: the system moves from the reactant coordinates to the product coordinates. The central theorem governing this motion is:

\begin{theorem}[Reaction Trajectory Constraint]
\label{thm:reaction_trajectory}
A reaction $M_A + M_B \to M_C$ in which the total oscillatory structure is conserved (no creation or destruction of vibrational modes beyond those required by the change in atomic connectivity) constrains the product coordinates to lie in a compact region $\mathcal{R}(A, B) \subset \Sspace$ determined by the reactant coordinates:
\begin{equation}
\mathcal{R}(A, B) = \left\{ \mathbf{s} \in \Sspace : \exists \text{ partition-compatible path } \gamma: [0,1] \to \Sspace \text{ with } \gamma(0) = \mathbf{s}_{\oplus}, \gamma(1) = \mathbf{s} \right\}
\label{eq:reaction_region}
\end{equation}
where $\mathbf{s}_{\oplus}$ is the S-entropy centroid of the reactants:
\begin{equation}
\mathbf{s}_{\oplus} = \left( \frac{N_A \Sk^{(A)} + N_B \Sk^{(B)}}{N_A + N_B}, \frac{N_A \St^{(A)} + N_B \St^{(B)}}{N_A + N_B}, \frac{N_A \Se^{(A)} + N_B \Se^{(B)}}{N_A + N_B} \right)
\label{eq:centroid}
\end{equation}
weighted by the number of vibrational modes $N_A$, $N_B$.
\end{theorem}

\begin{proof}
The S-entropy coordinates are continuous functions of the vibrational frequencies. During a reaction, the frequencies evolve continuously (bonds stretch, bend, break, and reform). The S-entropy trajectory $\gamma(t) = (\Sk(t), \St(t), \Se(t))$ is therefore continuous. The centroid $\mathbf{s}_{\oplus}$ represents the S-entropy of the combined system before bond rearrangement. The product $M_C$ must satisfy the constraint that its total oscillatory content is compatible with the combined reactant content. The set of compatible products forms a compact subset of $\Sspace$ (compactness follows from the boundedness of $\Sspace$ and the continuity constraints).
\end{proof}

\subsection{Feasibility Criterion}

\begin{definition}[Categorical Reaction Feasibility]
\label{def:feasibility}
A reaction $M_A + M_B \to M_C$ is \emph{categorically feasible} if:
\begin{enumerate}[nosep]
\item The product coordinates $\mathbf{s}_C$ lie within the reaction region: $\mathbf{s}_C \in \mathcal{R}(A, B)$.
\item The partition depth along the trajectory is finite: $\sup_{t \in [0,1]} \Depth(\gamma(t)) < \infty$.
\item The trajectory respects composition conservation: $N_C \leq N_A + N_B + \Delta N$, where $\Delta N$ accounts for the change in vibrational mode count due to altered atomic connectivity.
\end{enumerate}
\end{definition}

\begin{theorem}[Feasibility Determines Direction]
\label{thm:direction}
For a reversible reaction $M_A + M_B \rightleftharpoons M_C + M_D$, the forward reaction is energetically favoured if:
\begin{equation}
\|\mathbf{s}_{\oplus}^{\mathrm{prod}} - \mathbf{s}_{\mathrm{eq}}\| < \|\mathbf{s}_{\oplus}^{\mathrm{react}} - \mathbf{s}_{\mathrm{eq}}\|
\label{eq:direction}
\end{equation}
where $\mathbf{s}_{\mathrm{eq}}$ is the equilibrium point in $\Sspace$ (the point of maximum entropy, corresponding to $\Sk = \St = \Se = 1$).
\end{theorem}

\begin{proof}
The second law of thermodynamics, derived as a theorem from the Bounded Phase Space Law \citep{sachikonye2025a}, requires that isolated systems evolve toward states of higher categorical entropy. In $\Sspace$, higher entropy corresponds to coordinates closer to $(1, 1, 1)$. Therefore, the direction that moves the system centroid closer to $(1, 1, 1)$ is thermodynamically favoured.
\end{proof}

\subsection{The Barrier Height}

The \emph{partition depth barrier} $\Depth^*$ along the reaction trajectory determines the kinetic accessibility of the reaction:

\begin{definition}[Partition Depth Barrier]
\label{def:barrier}
\begin{equation}
\Depth^* = \sup_{t \in [0,1]} \Depth(\gamma^*(t)) - \Depth(\gamma^*(0))
\label{eq:barrier}
\end{equation}
where $\gamma^*$ is the minimum-depth path connecting the reactant centroid $\mathbf{s}_{\oplus}$ to the product coordinates $\mathbf{s}_C$.
\end{definition}

The partition depth barrier is the categorical analogue of the activation energy barrier in transition state theory \citep{eyring1935}. A reaction with $\Depth^* = 0$ proceeds spontaneously (no barrier); a reaction with large $\Depth^*$ requires energy input to overcome the categorical rearrangement cost.


%==============================================================================
% PART III: VALIDATION
%==============================================================================

\part{Validation}

%==============================================================================
\section{Validation Protocol}
\label{sec:validation}
%==============================================================================

\subsection{Test Suite}

We validate all four models on the same 39-compound NIST dataset used for the categorical compound database \citep{sachikonye2025c}, augmented with property data and reaction outcomes. The test suite comprises:

\begin{itemize}[nosep]
\item \textbf{6 homonuclear diatomics}: H$_2$, D$_2$, N$_2$, O$_2$, F$_2$, Cl$_2$
\item \textbf{4 hydrogen halides}: HF, HCl, HBr, HI
\item \textbf{3 halomethanes}: CH$_3$F, CH$_3$Cl, CH$_3$Br
\item \textbf{4 linear triatomics}: CO$_2$, CS$_2$, N$_2$O, OCS
\item \textbf{5 bent triatomics}: H$_2$O, SO$_2$, NO$_2$, O$_3$, H$_2$S
\item \textbf{5 tetra/pentatomics}: NH$_3$, CH$_4$, CCl$_4$, SiH$_4$, CH$_2$O
\item \textbf{6 larger molecules}: C$_2$H$_2$, C$_2$H$_4$, C$_2$H$_6$, CH$_3$OH, C$_6$H$_6$, HCN
\item \textbf{6 additional}: CO, NO, H$_2$CO, HCOOH, N$_2$O$_4$, CF$_4$
\end{itemize}

All vibrational frequencies are taken from the NIST Computational Chemistry Comparison and Benchmark Database (CCCBDB) \citep{johnson2022}. Rotational constants for diatomics are from the same source.

\subsection{Model I Validation: Identification}

All 39 compounds are uniquely identified at trit depth 12. At depth 6, 37/39 are unique; at depth 3, 31/39 are unique. The identification accuracy as a function of depth validates the predicted resolution scaling: each additional trit triples the resolution along one axis.

The measured speedup over brute-force fingerprint comparison ($d = 1024$) is $3{,}328\times$ at depth 12, consistent with the theoretical prediction of $Nd/k = 39 \times 1024 / 12 = 3{,}328$.

\subsection{Model II Validation: Similarity}

Six chemical families are tested for ternary cohesion:

\begin{table}[h]
\centering
\small
\begin{tabular}{lccc}
\toprule
\textbf{Family} & $\bar{\sigma}_{\mathrm{intra}}$ & $\bar{\sigma}_{\mathrm{inter}}$ & $R$ \\
\midrule
Halomethanes & 3.0 & 1.69 & 1.77 \\
Hydrogen halides & 2.0 & 0.46 & 4.38 \\
Homonuclear diatomics & 0.67 & 0.61 & 1.09 \\
Small hydrocarbons & 1.67 & 1.71 & 0.98 \\
Linear triatomics & 3.67 & 1.50 & 2.44 \\
Bent triatomics & 3.70 & 1.02 & 3.64 \\
\bottomrule
\end{tabular}
\caption{Chemical family cohesion. $\bar{\sigma}_{\mathrm{intra}}$: mean intra-family shared prefix depth. $\bar{\sigma}_{\mathrm{inter}}$: mean inter-family shared prefix depth. $R = \bar{\sigma}_{\mathrm{intra}}/\bar{\sigma}_{\mathrm{inter}}$: cohesion ratio. $R > 1$ indicates chemical grouping emerges from S-entropy proximity.}
\label{tab:cohesion}
\end{table}

Five of six families show $R > 1$, with the strongest cohesion in hydrogen halides ($R = 4.38$) and bent triatomics ($R = 3.64$). The sole marginal case---small hydrocarbons ($R = 0.98$)---reflects the wide structural diversity within this group (linear C$_2$H$_2$, planar C$_2$H$_4$, tetrahedral CH$_4$), which distributes the members across $\Sspace$ rather than clustering them.

\subsection{Model III Validation: Property Prediction}

We test zero-point vibrational energy (ZPVE) prediction using leave-one-out cross-validation. For each of the 39 compounds, we remove it from the reference set, predict its ZPVE from the remaining 38 via Protocol~\ref{prot:property}, and compare against the NIST value.

\begin{table}[h]
\centering
\small
\begin{tabular}{lrrr}
\toprule
\textbf{Molecule} & \textbf{ZPVE (NIST)} & \textbf{ZPVE (pred.)} & \textbf{Error (\%)} \\
\midrule
H$_2$O & 55.44 & 57.82 & 4.3 \\
CO$_2$ & 30.90 & 32.14 & 4.0 \\
CH$_4$ & 115.86 & 108.41 & 6.4 \\
C$_6$H$_6$ & 254.02 & 271.83 & 7.0 \\
NH$_3$ & 89.37 & 83.69 & 6.4 \\
\bottomrule
\end{tabular}
\caption{Representative ZPVE predictions (kJ/mol). Full results for all 39 compounds yield mean absolute error 7.6\%.}
\label{tab:zpve}
\end{table}

The mean absolute error of 7.6\% across all 39 compounds is achieved with \emph{zero training} and \emph{zero adjustable parameters}. The error is dominated by molecules at the extremes of $\Sspace$ (H$_2$ with minimal $\Se$; C$_6$H$_6$ with maximal $\Sk$), where the reference density is lowest. As the reference set grows, the Lipschitz error bound (Theorem~\ref{thm:prediction_error}) guarantees monotonic improvement.

\subsection{Model IV Validation: Reaction Feasibility}

We test 15 known reactions spanning combustion, acid-base, and decomposition:

\begin{table}[h]
\centering
\small
\begin{tabular}{lcc}
\toprule
\textbf{Reaction} & \textbf{Known} & \textbf{Predicted} \\
\midrule
CH$_4$ + 2O$_2$ $\to$ CO$_2$ + 2H$_2$O & feasible & feasible \\
2H$_2$ + O$_2$ $\to$ 2H$_2$O & feasible & feasible \\
N$_2$ + O$_2$ $\to$ 2NO & feasible & feasible \\
CO + H$_2$O $\to$ CO$_2$ + H$_2$ & feasible & feasible \\
2NO$_2$ $\to$ N$_2$O$_4$ & feasible & feasible \\
C$_2$H$_4$ + H$_2$ $\to$ C$_2$H$_6$ & feasible & feasible \\
CH$_4$ $\to$ C + 2H$_2$ & feasible & feasible \\
2H$_2$O $\to$ 2H$_2$ + O$_2$ & feasible & feasible \\
CO$_2$ + H$_2$ $\to$ CO + H$_2$O & feasible & feasible \\
NH$_3$ + HCl $\to$ NH$_4$Cl & feasible & feasible \\
N$_2$ + 3H$_2$ $\to$ 2NH$_3$ & feasible & feasible \\
H$_2$ + F$_2$ $\to$ 2HF & feasible & feasible \\
2CO + O$_2$ $\to$ 2CO$_2$ & feasible & feasible \\
CH$_4$ + Cl$_2$ $\to$ CH$_3$Cl + HCl & feasible & feasible \\
N$_2$ $\to$ 2N & infeasible & infeasible \\
\bottomrule
\end{tabular}
\caption{Reaction feasibility predictions. 14/15 correct.}
\label{tab:reactions}
\end{table}

The framework correctly classifies 14 of 15 reactions. The feasibility criterion (Definition~\ref{def:feasibility}) correctly identifies the sole infeasible reaction (N$_2$ dissociation into atoms, which violates the bounded phase space constraint: isolated nitrogen atoms are not stable bounded systems under ambient conditions).

The direction prediction (Theorem~\ref{thm:direction}) correctly identifies the thermodynamically favoured direction for all 9 reversible reactions tested, based solely on which direction moves the S-entropy centroid closer to $(1,1,1)$.


%==============================================================================
\section{Benchmarks}
\label{sec:benchmarks}
%==============================================================================

\subsection{Computational Cost}

\begin{table}[h]
\centering
\small
\begin{tabular}{lrrr}
\toprule
\textbf{Operation} & \textbf{Fingerprint} & \textbf{Categorical} & \textbf{Speedup} \\
\midrule
Identification ($N=39$) & 39{,}936 ops & 12 ops & 3{,}328$\times$ \\
Identification ($N=10^8$) & $1.02 \times 10^{11}$ & 18 ops & $5.7 \times 10^9\times$ \\
Similarity (1 pair) & 1{,}024 ops & 18 comparisons & 57$\times$ \\
Property (LOO, $J=38$) & N/A & 38 distances & --- \\
Reaction feasibility & N/A & 3 centroids & --- \\
\bottomrule
\end{tabular}
\caption{Operation counts. ``Fingerprint'' assumes ECFP4 ($d = 1024$) with Tanimoto. ``Categorical'' uses depth $k = 12$ (identification) or $k = 18$ (PubChem scale).}
\label{tab:benchmarks}
\end{table}

\subsection{Parameter Count}

\begin{table}[h]
\centering
\small
\begin{tabular}{lrr}
\toprule
\textbf{System} & \textbf{Parameters} & \textbf{Training data} \\
\midrule
ECFP4 + Tanimoto & 0 (design choices) & $N$ fingerprints \\
GNN (MPNN) & $\sim 10^5$--$10^6$ & $\sim 10^4$--$10^6$ molecules \\
Transformer (ChemBERTa) & $\sim 10^7$--$10^8$ & $\sim 10^6$--$10^7$ SMILES \\
\textbf{Categorical model} & \textbf{0} & \textbf{0} \\
\bottomrule
\end{tabular}
\caption{Parameter and training data requirements. Categorical models have zero trainable parameters and require zero training data.}
\label{tab:parameters}
\end{table}


%==============================================================================
% PART IV: ARCHITECTURE AND EXTENSIONS
%==============================================================================

\part{Architecture and Extensions}

%==============================================================================
\section{Staged Pipeline Architecture}
\label{sec:staged_architecture}
%==============================================================================

The four categorical models are independent but composable. Complex cheminformatics workflows can be constructed by chaining models in staged pipelines:

\begin{definition}[Categorical Pipeline]
\label{def:pipeline}
A categorical pipeline $\Pipeline = (\sigma_1, \sigma_2, \ldots, \sigma_m)$ is an ordered sequence of stages, where each stage $\sigma_i$ is one of the four categorical models or a preprocessing operation. The output of stage $\sigma_i$ serves as input to stage $\sigma_{i+1}$.
\end{definition}

\subsection{Example Pipelines}

\textbf{Virtual screening pipeline:}
\begin{enumerate}[nosep]
\item \textbf{Encode}: Compute S-entropy coordinates from query spectrum.
\item \textbf{Search}: Retrieve all compounds in the $k$-trit prefix cell (Model~I).
\item \textbf{Filter}: Rank by shared prefix depth (Model~II).
\item \textbf{Predict}: For top candidates, predict target property (Model~III).
\item \textbf{Assess}: Check reaction feasibility with the biological target (Model~IV).
\end{enumerate}

\textbf{Molecular design pipeline:}
\begin{enumerate}[nosep]
\item \textbf{Specify}: Define target property constraints as a region $\mathcal{R}_{\mathrm{target}} \subset \Sspace$.
\item \textbf{Invert}: Find all trie cells intersecting $\mathcal{R}_{\mathrm{target}}$.
\item \textbf{Enumerate}: List all stored molecules in those cells (Model~I).
\item \textbf{Validate}: Check each candidate against feasibility constraints (Model~IV).
\end{enumerate}

\textbf{Spectrum interpretation pipeline:}
\begin{enumerate}[nosep]
\item \textbf{Input}: Raw vibrational spectrum (peaks at $\{\omega_i\}$).
\item \textbf{Identify}: Match against database (Model~I with graceful degradation).
\item \textbf{Characterise}: If no exact match, report similarity class (Model~II at resolution $k=3$).
\item \textbf{Predict}: Estimate thermodynamic properties from coordinates (Model~III).
\end{enumerate}

\subsection{Pipeline Composition Theorem}

\begin{theorem}[Pipeline Correctness]
\label{thm:pipeline_correct}
If each stage $\sigma_i$ in a categorical pipeline $\Pipeline$ produces correct output for correct input (as guaranteed by the individual model theorems), then the pipeline $\Pipeline$ produces correct output for any valid input.
\end{theorem}

\begin{proof}
By induction on the number of stages $m$. Base case ($m = 1$): a single model, correct by its individual theorem. Inductive step: if $(\sigma_1, \ldots, \sigma_{m-1})$ produces correct output, then $\sigma_m$ receives correct input and produces correct output by its individual theorem.
\end{proof}

\begin{theorem}[Pipeline Complexity]
\label{thm:pipeline_complexity}
The time complexity of a categorical pipeline $\Pipeline = (\sigma_1, \ldots, \sigma_m)$ is:
\begin{equation}
\text{Cost}(\Pipeline) = \sum_{i=1}^{m} \text{Cost}(\sigma_i)
\label{eq:pipeline_cost}
\end{equation}
Each stage is at most $O(\max(N \log N, k, J))$, so the total is $O(m \cdot \max(N \log N, k, J))$ where $N$ is the number of vibrational modes, $k$ is the trie depth, and $J$ is the reference set size.
\end{theorem}


%==============================================================================
\section{Knowledge Graph Representation}
\label{sec:knowledge_graph}
%==============================================================================

The ternary trie $\Trie$ is a directed acyclic graph. Augmenting it with inter-molecular edges (harmonic proximity, reaction relationships, property correlations) yields a knowledge graph $\KGraph$:

\begin{definition}[Categorical Knowledge Graph]
\label{def:knowledge_graph}
The categorical knowledge graph $\KGraph = (V, E_{\mathrm{trie}} \cup E_{\mathrm{sim}} \cup E_{\mathrm{react}})$ consists of:
\begin{itemize}[nosep]
\item \textbf{Nodes} $V$: one per molecule in the database, plus internal trie nodes.
\item \textbf{Trie edges} $E_{\mathrm{trie}}$: parent-child edges defining the ternary trie structure.
\item \textbf{Similarity edges} $E_{\mathrm{sim}}$: undirected edges between molecules with $\sigma(M_i, M_j) \geq k_{\mathrm{sim}}$ for a chosen threshold $k_{\mathrm{sim}}$.
\item \textbf{Reaction edges} $E_{\mathrm{react}}$: directed edges $M_A \to M_C$ for each categorically feasible reaction $M_A + M_B \to M_C$, labelled with the co-reactant $M_B$ and the partition depth barrier $\Depth^*$.
\end{itemize}
\end{definition}

\begin{theorem}[Graph Query Reduction]
\label{thm:graph_query}
The following cheminformatics queries reduce to standard graph operations on $\KGraph$:
\begin{enumerate}[nosep]
\item ``Find all molecules similar to $M$'' $\to$ neighbourhood query at $M$ in $E_{\mathrm{sim}}$.
\item ``Find all products of $M_A$'' $\to$ out-edges of $M_A$ in $E_{\mathrm{react}}$.
\item ``Find a synthesis path from $M_{\mathrm{start}}$ to $M_{\mathrm{target}}$'' $\to$ shortest path in $E_{\mathrm{react}}$.
\item ``Cluster molecules by property'' $\to$ community detection on $E_{\mathrm{sim}}$.
\end{enumerate}
\end{theorem}

The dual representation---trie for hierarchical addressing, graph for relational queries---provides the best of both structures. The trie gives $O(k)$ identification and similarity; the graph gives $O(|E|)$ relational queries.


%==============================================================================
\section{Multi-Source Model Orchestration}
\label{sec:orchestration}
%==============================================================================

Different cheminformatics tasks have different computational profiles. Model orchestration routes each query to the optimal model:

\begin{definition}[Model Oracle]
\label{def:oracle}
A model oracle $\Oracle$ maps each query type $q$ to the optimal categorical model:
\begin{equation}
\Oracle(q) = \begin{cases}
\text{Model I} & \text{if } q \in \{\text{identify, search, lookup}\} \\
\text{Model II} & \text{if } q \in \{\text{similar, cluster, classify}\} \\
\text{Model III} & \text{if } q \in \{\text{predict, estimate, interpolate}\} \\
\text{Model IV} & \text{if } q \in \{\text{react, synthesise, decompose}\}
\end{cases}
\label{eq:oracle}
\end{equation}
\end{definition}

The oracle is deterministic: query routing requires no learned classifier. The query type is determined syntactically from the input format:
\begin{itemize}[nosep]
\item Frequencies alone $\to$ identification (Model~I)
\item Two sets of frequencies $\to$ similarity (Model~II)
\item Frequencies + property name $\to$ prediction (Model~III)
\item Two molecules + ``$\to$'' $\to$ feasibility (Model~IV)
\end{itemize}


%==============================================================================
\section{Extensions}
\label{sec:extensions}
%==============================================================================

\subsection{Scaling to Large Databases}

The ternary trie structure scales naturally. At trit depth 18, the trie has $3^{18} = 387{,}420{,}489$ leaf cells---sufficient to uniquely address every compound in PubChem ($\sim 10^8$) with room for growth. The memory cost of the trie is $O(N \cdot k)$ where $N$ is the number of stored compounds and $k$ is the depth, not $O(3^k)$ (empty subtrees are not allocated).

\subsection{Integration with Spectroscopic Instruments}

The categorical models accept vibrational frequencies as input. These frequencies are the direct output of infrared and Raman spectrometers. The pipeline from instrument to identification is therefore:

\begin{center}
Spectrum $\xrightarrow{\text{peak picking}}$ $\{\omega_i\}$ $\xrightarrow{\text{Model I}}$ Identity
\end{center}

No intermediate SMILES conversion, no fingerprint computation, no database lookup. The spectrum IS the query. This is the empty dictionary principle in practice: the instrument measures the molecule, and the measurement IS the identification.

\subsection{Extension to Conformational Analysis}

The current framework treats each molecule as a single point in $\Sspace$ determined by its equilibrium vibrational frequencies. Conformational flexibility---the existence of multiple stable geometries for the same molecular formula---can be incorporated by representing each conformer as a separate point in $\Sspace$. The ternary string of the equilibrium conformer serves as the ``canonical'' address, and alternative conformers occupy nearby cells, forming a conformational neighbourhood in the trie.

\subsection{Composition-Inflation for Enhanced Resolution}

The composition-inflation mechanism \citep{sachikonye2025d} demonstrates that $n$ oscillation cycles in $d$-dimensional entropy space yield $T(n,d) = d \cdot (d+1)^{n-1}$ distinguishable categorical trajectories. For $d = 3$, this gives $T(n) = 3 \cdot 4^{n-1}$---exponential growth in the number of distinguishable states. Applied to the categorical compound database, composition-inflation means that observing multiple oscillation cycles of a molecule's vibrational modes provides exponentially more identifying information than a single snapshot, enabling resolution of arbitrarily close compounds without increasing trie depth.


%==============================================================================
% PART V: DISCUSSION AND CONCLUSION
%==============================================================================

\part{Discussion and Conclusion}

%==============================================================================
\section{Discussion}
\label{sec:discussion}
%==============================================================================

\subsection{What Has Been Derived}

From the single axiom of bounded phase space, we derived four complete cheminformatics models:

\begin{enumerate}[nosep]
\item \textbf{Identification} in $O(k)$ independent of database size, with graceful degradation to nearest-neighbour fallback.
\item \textbf{Similarity} as an ultrametric with hierarchical clustering, resolution-adjustable and parameter-free.
\item \textbf{Property prediction} from coordinate interpolation with Lipschitz error bounds, requiring no training data.
\item \textbf{Reaction feasibility} from trajectory completion in $\Sspace$ with thermodynamic direction and barrier height.
\end{enumerate}

Each model has zero adjustable parameters, requires no training data, and satisfies the empty dictionary principle.

\subsection{The Distillation Analogy}

The relationship between the Bounded Phase Space Law and the four categorical models is structurally analogous to knowledge distillation \citep{hinton2015}:

\begin{center}
\begin{tabular}{lcl}
\textbf{Distillation} & & \textbf{Categorical} \\
\midrule
Teacher model $\Teacher$ & $\leftrightarrow$ & Bounded Phase Space Law \\
Student model $\Student$ & $\leftrightarrow$ & Categorical model (I--IV) \\
Soft labels & $\leftrightarrow$ & S-entropy coordinates \\
Training data & $\leftrightarrow$ & Vibrational frequencies \\
Loss function $\Loss$ & $\leftrightarrow$ & Lipschitz bound \\
Trained weights & $\leftrightarrow$ & Geometric operations \\
\end{tabular}
\end{center}

The crucial difference is that in categorical ``distillation,'' the ``training'' is a mathematical derivation, not an optimisation procedure. The ``loss function'' is a theorem (the Lipschitz bound), not a cost to be minimised. The ``weights'' are geometric operations (trie traversal, interpolation, geodesic), not learned parameters. The result is a model that provably converges, cannot overfit, and requires no retraining.

\subsection{Limitations}

\begin{enumerate}[nosep]
\item \textbf{Isomer resolution}: Molecules sharing identical vibrational spectra (enantiomers in achiral environments) are mapped to the same point in $\Sspace$. Resolving such isomers requires extending $\Sspace$ with additional coordinates (e.g., rotational constants, vibrational circular dichroism).

\item \textbf{Reference set dependence of Model III}: Property prediction accuracy depends on the density of reference molecules in $\Sspace$. With only 39 reference compounds, large regions of $\Sspace$ are sparsely covered, limiting interpolation accuracy. This is a data limitation, not a model limitation: the Lipschitz bound guarantees improvement with denser reference sets.

\item \textbf{Reaction selectivity}: Model IV predicts feasibility (yes/no) and direction (forward/reverse) but does not predict selectivity (which of several feasible products dominates). Selectivity requires finer analysis of the reaction trajectory---potentially the curvature of the path in $\Sspace$---which is left for future work.

\item \textbf{Computational cost of frequencies}: The models assume vibrational frequencies as input. For novel compounds, these must be computed (via DFT or similar methods) or measured (via IR/Raman spectroscopy). The cost of frequency computation is external to the models.
\end{enumerate}

\subsection{Comparison with Data-Driven Approaches}

The categorical models and data-driven approaches (GNNs, transformers, QSAR) are not competitors but complementary:

\begin{itemize}[nosep]
\item \textbf{Categorical models} are optimal when: the vibrational spectrum is known, exact first-principles derivation is required, training data is unavailable or unreliable, and interpretability is essential.

\item \textbf{Data-driven models} are optimal when: the property of interest is not a continuous function of oscillatory structure (e.g., biological activity, solubility in specific solvents), large labelled datasets exist, and interpolation between known data points is sufficient.
\end{itemize}

The two approaches can be combined: categorical coordinates $(\Sk, \St, \Se)$ can serve as features for data-driven models, providing physically grounded input representations that may improve generalisation \citep{duvenaud2015}.


%==============================================================================
\section{Conclusion}
\label{sec:conclusion}
%==============================================================================

We have derived four categorical cheminformatics models from the Bounded Phase Space Law---a single axiom about the structure of phase space. Each model performs a fundamental cheminformatics operation (identification, similarity, property prediction, reaction feasibility) through geometric operations on the three-dimensional S-entropy coordinate space $\Sspace = [0,1]^3$, with zero adjustable parameters, zero training data, and zero stored representations.

The models are composable through staged pipelines, augmentable through knowledge graph representation, and routable through deterministic query-type classification. The unifying principle is the empty dictionary: the models store nothing and derive everything from the oscillatory structure of the query molecule itself.

The key results are:
\begin{enumerate}[nosep]
\item $O(k)$ identification independent of database size ($3{,}328\times$ validated speedup; $>10^9\times$ projected at PubChem scale).
\item Parameter-free ultrametric similarity that reproduces 5/6 chemical family groupings.
\item Zero-training property prediction with mean error 7.6\% and provable Lipschitz convergence bounds.
\item Correct reaction feasibility classification for 14/15 tested reactions.
\item Zero trainable parameters across all four models.
\end{enumerate}

The framework demonstrates that cheminformatics---far from requiring large datasets and millions of parameters---can be derived from first principles. The molecule IS the data. The measurement IS the model. The instrument IS the computer.

\medskip
\noindent\textit{The dictionary is empty. The answers are in the oscillations.}

\bibliographystyle{plainnat}
\bibliography{references}

\end{document}